% 子文件：G_连线综合_小学
% 本文件包含 6 个 section



%% ==============================
%% 连线题：数字中“5”的意义
%% ==============================
\newpage

\section{解答：连线题——“5”表示什么}

\vspace{0.4cm}

\noindent\textbf{2. 下面各数中的“5”表示什么？连一连。}

\vspace{0.8cm}

\begin{center}
\begin{tikzpicture}[>=Stealth, scale=1.0] % 整体思路：先把上排数字和下排“5表示的意义”分别放成两行节点，再利用节点锚点 south、north 画对应连线，形成清晰的匹配图
  % 顶部数字节 (T1 到 T5)
  \node (T1) at (0, 2) {\large $0.405$}; % 建立名为 T1 的节点；at (0,2) 指定位置；\large 放大字号；内容是数字 0.405
  \node (T2) at (3, 2) {\large $50.03$}; % 顶排第二个数字节点，横向右移到 x=3
  \node (T3) at (6, 2) {\large $12.85$}; % 顶排第三个数字节点，保持同一高度 y=2
  \node (T4) at (10, 2) {\large $0.577$}; % 顶排第四个数字节点，x=10 使间距更舒展
  \node (T5) at (14, 2) {\large $105.3$}; % 顶排第五个数字节点，放在最右侧

  % 底部文字节点 (B1 到 B5)
  \node (B1) at (0, 0) {\large 5个十}; % 建立底排第一个含义节点；名称 B1 便于后面连线调用
  \node (B2) at (3, 0) {\large 5个一}; % 底排第二个含义节点
  \node (B3) at (6, 0) {\large 5个十分之一}; % 底排第三个含义节点
  \node (B4) at (10, 0) {\large 5个百分之一}; % 底排第四个含义节点
  \node (B5) at (14, 0) {\large 5个千分之一}; % 底排第五个含义节点

  % 连线 (红色实线，从南到北)
  % 0.405 中 5 在千分位 -> 5个千分之一 (T1 -> B5)
  \draw[red!70!black, line width=1.5pt] (T1.south) -- (B5.north); % 从上方节点 T1 的下锚点 south 连到底部节点 B5 的上锚点 north；用红色粗线突出配对关系
  
  % 50.03 中 5 在十位 -> 5个十 (T2 -> B1)
  \draw[red!70!black, line width=1.5pt] (T2.south) -- (B1.north); % 从 T2 向 B1 画匹配线，表示 50.03 里的 5 代表 5个十

  % 12.85 中 5 在百分位 -> 5个百分之一 (T3 -> B4)
  \draw[red!70!black, line width=1.5pt] (T3.south) -- (B4.north); % 从 T3 连到 B4，说明 12.85 中的 5 在百分位

  % 0.577 中 5 在十分位 -> 5个十分之一 (T4 -> B3)
  \draw[red!70!black, line width=1.5pt] (T4.south) -- (B3.north); % 从 T4 连到 B3，表示 0.577 中的 5 在十分位

  % 105.3 中 5 在个位 -> 5个一 (T5 -> B2)
  \draw[red!70!black, line width=1.5pt] (T5.south) -- (B2.north); % 从 T5 连到 B2，表示 105.3 中的 5 在个位

\end{tikzpicture} % 结束“数位意义连线图”
\end{center}

\vspace{0.6cm}

{\small\color{gray}
\textbf{解析：}\\
$0.405$ 的“$5$”在千分位，表示 $5$ 个千分之一；\\
$50.03$ 的“$5$”在十位，表示 $5$ 个十；\\
$12.85$ 的“$5$”在百分位，表示 $5$ 个百分之一；\\
$0.577$ 的“$5$”在十分位，表示 $5$ 个十分之一；\\
$105.3$ 的“$5$”在个位，表示 $5$ 个一。
}

\vspace{0.6cm}

{\small\color{blue!70!black}
\textbf{绘图基本原理与步骤：}\\
\textbf{1. 节点对称布局}：在 $y=2$ 和 $y=0$ 两条平行线上分别摆放上排数字节点（T1--T5）和下排含义节点（B1--B5），形成匹配布局。\\
\textbf{2. 锚点连线}：利用 \texttt{T.south} 和 \texttt{B.north} 调用南北锚点连线，交叉自然穿梭不因锚点偏移而错位。\\
\textbf{3. 视觉配对强调}：使用红色粗线 \texttt{red!70!black, line width=1.5pt} 统一突出匹配连线。
}

%% ==============================
%% 画图解题：相向而行与相背而行
%% ==============================
\newpage



%% ==============================
%% 大数的认识：各国陆地面积
%% ==============================
\newpage

\section{解答：世界上陆地面积最大的四个国家}

\vspace{0.4cm}

\noindent\textbf{4. 世界上陆地面积最大的四个国家的面积如下表：}

\vspace{0.4cm}

\begin{center}
\renewcommand{\arraystretch}{1.6}
\setlength{\tabcolsep}{15pt}
\begin{tabular}{|c|c|}
\hline
中\quad 国 & 九百六十万平方千米 \\
\hline
俄罗斯 & 一千七百零九万八千二百平方千米 \\
\hline
美\quad 国 & 九百三十七万平方千米 \\
\hline
加拿大 & 九百九十八万平方千米 \\
\hline
\end{tabular}
\end{center}

\vspace{0.6cm}

\noindent\textbf{（1）填写下表。}

\vspace{0.4cm}

\begin{center}
\renewcommand{\arraystretch}{1.6}
\setlength{\tabcolsep}{12pt}
\begin{tabular}{|c|c|c|}
\hline
国\quad 家 & 面积/平方千米 & 面积/万平方千米 \\
\hline
中\quad 国 & \textcolor{red}{9600000} & \textcolor{red}{960} \\
\hline
俄罗斯 & \textcolor{red}{17098200} & \textcolor{red}{1709.82} \\
\hline
美\quad 国 & \textcolor{red}{9370000} & \textcolor{red}{937} \\
\hline
加拿大 & \textcolor{red}{9980000} & \textcolor{red}{998} \\
\hline
\end{tabular}
\end{center}

\vspace{0.6cm}
{\small\color{gray}
\textbf{解析：}\\
首先写出各面积的阿拉伯数字形式：\\
九百六十万 $\rightarrow$ 9600000\\
一千七百零九万八千二百 $\rightarrow$ 17098200\\
九百三十七万 $\rightarrow$ 9370000\\
九百九十八万 $\rightarrow$ 9980000\\
将其改写为以“万”为单位的数，只需将小数点向左移动 4 位：\\
$9600000 = 960$ 万；
$17098200 = 1709.82$ 万；
$9370000 = 937$ 万；
$9980000 = 998$ 万。
}

\vspace{0.8cm}

%% ==============================
%% 分数的初步认识：涂色表示分数
%% ==============================
\newpage




%% ==============================
%% 第2题：体育用品商店购买问题——单价、数量的关系
%% ==============================

\section{解答：体育用品商店购买问题}

\vspace{0.4cm}

\noindent\textbf{2. 一个体育用品商店，每个排球 30 元，每个足球 50 元。李老师带的钱正好能买 12 个篮球或 10 个排球。（先根据问题填写表格，再列式解答）}

\vspace{0.6cm}

% 填写表格
\begin{center}
\renewcommand{\arraystretch}{1.8}
\begin{tabular}{|c|c|c|}
\noalign{\hrule height 1pt}
\makebox[3cm][c]{品名} & \makebox[4cm][c]{单价 / (元 / 个)} & \makebox[3cm][c]{数量 / 个} \\
\hline
排球 & \textcolor{blue}{\textbf{30}} & \textcolor{blue}{\textbf{10}} \\
\hline
篮球 & \textcolor{blue}{\textbf{?}} & \textcolor{blue}{\textbf{12}} \\
\hline
足球 & \textcolor{blue}{\textbf{50}} & \textcolor{blue}{\textbf{?}} \\
\noalign{\hrule height 1pt}
\end{tabular}
\end{center}

\vspace{0.8cm}

\noindent\textbf{（1）每个篮球多少元？}

\vspace{0.2cm}

{\small\color{gray}
\textbf{解析：}\\
第一步，先求出李老师带的总钱数：排球单价 $\times$ 排球数量 $= 30 \times 10 = 300$（元）。\\
第二步，由于这笔钱正好也能买 12 个篮球，所以篮球的单价为：总钱数 $\div$ 篮球数量 $= 300 \div 12 = 25$（元）。
}

\vspace{0.2cm}
\noindent\textbf{综合算式解答：} \textbf{\textcolor{blue}{$\boldsymbol{30 \times 10 \div 12 = 25}$}} \textbf{（元）} \\
\noindent\textbf{答：}每个篮球 \textbf{\textcolor{blue}{25}} 元。

\vspace{0.6cm}

\noindent\textbf{（2）李老师带的钱能买多少个足球？}

\vspace{0.2cm}

{\small\color{gray}
\textbf{解析：}\\
同样利用前面求出的总钱数 300 元。足球的单价是 50 元/个，那么这笔钱能买足球的数量为：总钱数 $\div$ 足球单价 $= 300 \div 50 = 6$（个）。
}

\vspace{0.2cm}
\noindent\textbf{综合算式解答：} \textbf{\textcolor{blue}{$\boldsymbol{30 \times 10 \div 50 = 6}$}} \textbf{（个）} \\
\noindent\textbf{答：}李老师带的钱能买 \textbf{\textcolor{blue}{6}} 个足球。

\vspace{0.6cm}

\noindent\textbf{（3）如果买 12 个篮球和 8 个足球，一共要多少元？}

\vspace{0.2cm}

{\small\color{gray}
\textbf{解析：}\\
买 12 个篮球所需的钱就是李老师原先的总钱数（$300$ 元），或者用刚才算出的篮球单价计算：$25 \times 12 = 300$（元）。\\
买 8 个足球需要的钱为：$50 \times 8 = 400$（元）。\\
两者加起来一共要：$300 + 400 = 700$（元）。
}

\vspace{0.2cm}
\noindent\textbf{综合算式解答：} \textbf{\textcolor{blue}{$\boldsymbol{25 \times 12 + 50 \times 8 = 700}$}} \textbf{（元）} \ \textit{或直接写为：} \textbf{\textcolor{blue}{$\boldsymbol{300 + 50 \times 8 = 700}$}} \textbf{（元）}\\
\noindent\textbf{答：}一共要 \textbf{\textcolor{blue}{700}} 元。


\newpage



\newpage

\section{解答：假设法解决问题——象棋与跳棋数量}

\vspace{0.4cm}

\noindent\textbf{3. 学校有象棋、跳棋共 26 副，2 人下一副象棋，6 人下一副跳棋，恰好可以供 96 人进行活动。象棋、跳棋各有多少副？（先假设两种棋的数量，再调整找出答案）}

\vspace{0.8cm}

\begin{center}
\renewcommand{\arraystretch}{1.8}
% 表格绘制：4列，包含假设过程、计算过程及结果对比
\begin{tabular}{|c|c|l|c|}
\hline
\textbf{象棋数量/副} & \textbf{跳棋数量/副} & \multicolumn{1}{c|}{\textbf{下棋总人数}} & \textbf{和 96 人比较} \\ \hline
\textcolor{blue}{13} & \textcolor{blue}{13} & \textcolor{blue}{$13 \times 2 + 13 \times 6 = 104$} & \textcolor{blue}{多 8 人} \\ \hline
\textcolor{blue}{14} & \textcolor{blue}{12} & \textcolor{blue}{$14 \times 2 + 12 \times 6 = 100$} & \textcolor{blue}{多 4 人} \\ \hline
\textcolor{blue}{15} & \textcolor{blue}{11} & \textcolor{blue}{$15 \times 2 + 11 \times 6 = 96$} & \textcolor{blue}{相等} \\ \hline
\end{tabular}
\end{center}

\vspace{0.6cm}

{\small\color{gray}
\textbf{解析：}\\
1. **第一次假设**：假设象棋和跳棋各有 $13$ 副（$26 \div 2 = 13$）。计算总人数为 $13 \times 2 + 13 \times 6 = 104$（人），比 $96$ 人多了 $8$ 人。\\
2. **分析调整方向**：因为计算出的总人数偏多，且跳棋的人均负荷（6人）高于象棋（2人），因此需要减少跳棋数量。\\
3. **逐步调整**：每减少一副跳棋并增加一副象棋，总人数减少 $6 - 2 = 4$（人）。\\
4. **得出结论**：经过两次调整，当象棋为 $15$ 副、跳棋为 $11$ 副时，总人数刚好为 $96$ 人。
}

\vspace{0.4cm}
\noindent\textbf{解答：}\\
象棋有（ \textbf{\textcolor{blue}{15}} ）副，跳棋有（ \textbf{\textcolor{blue}{11}} ）副。

\vspace{0.6cm}

{\small\color{blue!70!black}
\textbf{绘图基本原理与步骤：}\\
\textbf{1. 列表尝试法（Guess and Check）}：通过列表有序展示不同假设下的计算结果，寻找变量变化与最终数值间的规律。\\
\textbf{2. 建立逻辑关联}：表格中每一行代表一次完整的业务逻辑模拟。第一、二列为独立变量，第三列为衍生计算量，第四列为反馈修正量。\\
\textbf{3. 数据一致性}：表格内所有数据项均使用 \texttt{textcolor} 进行着色处理，既突出了动态求解过程，也保持了与题干风格的视觉统一。\\
\textbf{4. 结构化呈现}：利用 \texttt{tabular} 的水平线与垂直线构筑严谨的逻辑网格，使抽象的代数试错过程变得直观易读。
}



\newpage



%% ==============================
%% 补充题：找出等式与方程并连一连
%% ==============================
\newpage

\section{解答：找出等式与方程并准确连线}

\vspace{0.6cm}

\begin{center}
\begin{tikzpicture}[>=Stealth, line join=round, line cap=round, scale=1.0]
  
  % 定点坐标参数 (四行题目，采用相等的纵向间距 1.5)
  \def\yA{4.5}
  \def\yB{3.0}
  \def\yC{1.5}
  \def\yD{0.0}

  % 定义全图共用的关键图元样式
  \tikzset{
    expr/.style={font=\Large, text=black},
    box/.style={draw=black, rectangle, minimum width=2.0cm, minimum height=1.0cm, font=\Large\bfseries, text=black, line width=0.8pt},
    conn/.style={draw=red, line width=1.4pt} % 还原参考答案阅卷风格的醒目红色加粗连线
  }

  % ===================================
  % 一、 结构排版布局：三列式设计
  % ===================================

  % 1. 左侧表达式列 (利用 anchor=east 使其右侧端点严格对其到 x=0 标线上)
  \node[expr, anchor=east] (L1) at (0, \yA) {$90 - x = 30$};
  \node[expr, anchor=east] (L2) at (0, \yB) {$80 \div 4 = 20$};
  \node[expr, anchor=east] (L3) at (0, \yC) {$x - 15$};
  \node[expr, anchor=east] (L4) at (0, \yD) {$7y = 63$};

  % 2. 居中分类框列 (等式框居于 1、2 行中间，方程框居于 3、4 行中间，完美匀称)
  \node[box] (M1) at (3.5, 3.75) {等式};
  \node[box] (M2) at (3.5, 0.75) {方程};

  % 3. 右侧表达式列 (利用 anchor=west 使其左侧端点严格对其到 x=7 标线上)
  \node[expr, anchor=west] (R1) at (7, \yA) {$20 + 30 = 50$};
  \node[expr, anchor=west] (R2) at (7, \yB) {$y + 17 = 38$};
  \node[expr, anchor=west] (R3) at (7, \yC) {$36 + x < 40$};
  \node[expr, anchor=west] (R4) at (7, \yD) {$54 \div x = 9$};

  % ===================================
  % 二、 逻辑自动连线
  % (利用直接指向节点名的方式，借助 TikZ 原生算法自动在方框轮廓线生成分散连接点)
  % ===================================

  % --- 【“等式”的连线群】 ---
  % 本质要求：只要含有等号“=”的式子统统属于等式
  % 从左侧连向等式：
  \draw[conn] (L1.east) -- (M1);
  \draw[conn] (L2.east) -- (M1);
  \draw[conn] (L4.east) -- (M1);
  % 从右侧连向等式：
  \draw[conn] (R1.west) -- (M1);
  \draw[conn] (R2.west) -- (M1);
  \draw[conn] (R4.west) -- (M1);

  % --- 【“方程”的连线群】 ---
  % 本质要求：在等式的基础上，含有未知数的式子才是方程
  % 从左侧连向方程：
  \draw[conn] (L1.east) -- (M2);
  \draw[conn] (L4.east) -- (M2);
  % 从右侧连向方程：
  \draw[conn] (R2.west) -- (M2);
  \draw[conn] (R4.west) -- (M2);

\end{tikzpicture}
\end{center}

\vspace{0.6cm}

{\small\color{gray}
\textbf{解题说明：}\\
1. \textbf{概念辨析：}表示相等关系的式子叫做\textbf{等式}，所有的等式中都必须含有等号（$=$）。含有未知数的等式叫做\textbf{方程}。因此，所有的方程一定是等式，但等式不一定是方程。\\
2. \textbf{左侧判断：}\\
\quad $90-x=30$：含有未知数 $x$ 且是等式，故既连“等式”也连“方程”；\\
\quad $80\div4=20$：是等式但不含未知数，只连“等式”；\\
\quad $x-15$：不含等号，既不是等式也不是方程；\\
\quad $7y=63$：含有未知数 $y$ 且是等式，既连“等式”也连“方程”。\\
3. \textbf{右侧判断：}\\
\quad $20+30=50$：是等式但不含未知数，只连“等式”；\\
\quad $y+17=38$：含未知数且是等式，既连“等式”也连“方程”；\\
\quad $36+x<40$：含有未知数但属于不等式（含 $<$），二者皆不可连；\\
\quad $54\div x=9$：含有未知数且是等式，既连“等式”也连“方程”。
}

\vspace{0.6cm}

{\small\color{blue!70!black}
\textbf{绘图基本原理与步骤：}\\
\textbf{1. 网格化布局节点：}通过精确定义常数如 \texttt{\textbackslash def\textbackslash yA\{4.5\}} 规范四行的绝对统一高度。左侧各题目节点赋予 \texttt{anchor=east} 属性锁定同一右侧端边界，右侧题目节点相应赋予 \texttt{anchor=west} 锁定同一左侧边。以此使整体横跨三竖排的结构极端工整齐平，还原经典教辅印刷品的空间视觉秩序感。\\
\textbf{2. 智能离散连线：}在 TikZ 算法中如果指定直接向节点框体中心点相连连线（即仅使用 \texttt{-- (M1)} 而不是精准定位锚点 \texttt{-- (M1.west)}），底层原生渲染器就会自动计算出每一条带有角度跨向中心的射线与方框边界线条的物理相交点，截断并锁止在那。这一代码小幅规避了所有线段极其死板地交汇同一点的生硬状态，呈现出仿佛手工在不同下笔位自然向中心聚拢划线的动态视觉效果。\\
\textbf{3. 醒目的红勾墨迹模拟：}利用预定义常量 \texttt{conn/.style=\{draw=red, line width=1.4pt\}} 为线系批量刷上大红粗线效果，跳出题目原有静止黑白色图层，打造最高效的查漏阅卷视角。
}



%% ==============================
%% 第27题（补充题）：方程特征判断
%% ==============================
\newpage

\section{解答：方程的意义与特征判断}

\vspace{0.6cm}

\textbf{\LARGE 12.} \textbf{下面式子中是方程的在后面方框里打“$\checkmark$”。}

\vspace{0.4cm}
\textbf{解：}

\textbf{（1）判断原理分析：}
“方程”在数学定义上必须同时满足两个核心条件：\textbf{① 必须是含有未知数的式子；② 必须是等式（含有“=”）}。
\begin{itemize}
    \item $x + 3.6 = 7$：含有未知数 $x$，且是等式。$\Rightarrow$ \textbf{是方程}。
    \item $4 \times 24 = 96$：虽然也是等式，但不含未知数字母（只是普通的算术等式）。$\Rightarrow$ \textbf{不是方程}。
    \item $a \times 2 > 2.4$：虽然含有未知数 $a$，但是包含“$>$”属于不等式。$\Rightarrow$ \textbf{不是方程}。
    \item $3 \div b = 1.5$：含有未知数 $b$，且属于等式架构。$\Rightarrow$ \textbf{是方程}。
    \item $5y + 2y$：虽然含有未知数 $y$，但这只是一个代数多项式，没有等号。$\Rightarrow$ \textbf{不是方程}。
    \item $2x + 3y = 9$：含有未知数 $x, y$，且是等式（标准的二元一次方程）。$\Rightarrow$ \textbf{是方程}。
\end{itemize}

\textbf{（2）选项勾画（原卷横向 $2 \times 3$ 三列方阵复位）：}

\vspace{0.4cm}
\begin{center}
\renewcommand{\arraystretch}{2.5}
% 启动自适应排版，将横向划分为三等分空间
\begin{tabular}{p{0.3\textwidth} p{0.3\textwidth} p{0.3\textwidth}}
{\Large $x + 3.6 = 7$} \checkedbox & {\Large $4 \times 24 = 96$} \emptybox & {\Large $a \times 2 > 2.4$} \emptybox \\
{\Large $3 \div b = 1.5$} \checkedbox & {\Large $5y + 2y$} \emptybox & {\Large $2x + 3y = 9$} \checkedbox \\
\end{tabular}
\end{center}

\vspace{0.4cm}
{\small\color{blue!70!black}
\textbf{画图及结构排版思路解构：}\\
\textbf{1. 方程指纹双因子安全重度校验算法：} 后台对这 6 个混合数学式直接执行了双层排雷侦测过滤。第一层拦截“未知数探测”：直接排除了没有变量字母作为主体的纯算术式（如 $4 \times 24 = 96$）；第二层拦截“等式平衡体特征验证（即寻找 \texttt{=} 符号）”：精准狙杀并剔除了不等式家族的叛逃者（$a \times 2 > 2.4$）和没有结项的纯粹多项式（$5y + 2y$）。最后沙盘中仅剩下 $x + 3.6 = 7$、$3 \div b = 1.5$ 和 $2x + 3y = 9$ 这三个完美携带着“未知数+等号”双螺旋基因的幸存项被系统赋予了代表胜利的红色对勾权限。\\
\textbf{2. 沿用悬浮自适应基线的超清打钩宏：} 制图部分我完美继承并共享了刚刚上两局锻造好的高精度 \texttt{\textbackslash{}emptybox} 与 \texttt{\textbackslash{}checkedbox} 引擎底层宏。这两个组件不仅尺寸锁死在了视觉最为舒适的 $0.4\mathrm{cm}$，还带有 \texttt{baseline=-0.15em} 的微缩下坠抓地力设计，这让它们哪怕在面对被刻意放大提拉的 \texttt{\textbackslash{}Large} 级别巨型数学公式时，其底盘基座依然能和前方数学符号的重心保持着工业级水准的一字水平防抖对齐！外部阵列这次彻底舍弃了死代码的物理绝对定距，转而采用以 \texttt{p\{0.3\textbackslash{}textwidth\}}（吃掉 $30\%$ 页宽）架构驱动的智能三等分液态框架，让这六组短式在庞大的整面白纸上依然能呈现出无比透气匀称的黄金分割布局！
}
%% ==============================
%% 第28题（补充题）：羊毛衫销售折线统计图
%% ==============================
\newpage


